\documentclass[14pt,aspectratio=169]{beamer}

% ── Theme ────────────────────────────────────────────────────────────
\usetheme{Madrid}
\usecolortheme{default}
\setbeamertemplate{navigation symbols}{}
\setbeamertemplate{footline}[frame number]

% ── Packages ─────────────────────────────────────────────────────────
\usepackage[ngerman]{babel}
\usepackage[utf8]{inputenc}
\usepackage[T1]{fontenc}
\usepackage{helvet}
\usepackage{graphicx}
\usepackage{tabularx}
\usepackage{booktabs}
\usepackage{xcolor}

% ── Farben ───────────────────────────────────────────────────────────
\definecolor{hsbl}{RGB}{0,82,136}
\setbeamercolor{structure}{fg=hsbl}
\setbeamercolor{title}{fg=white,bg=hsbl}
\setbeamercolor{frametitle}{fg=hsbl}

% ── Titel ────────────────────────────────────────────────────────────
\title{Ergonomia}
\subtitle{Ein ergonomisch gestalteter KI-Chatbot für das Immersive Reality Lab}
\author{Eigener Name | Matrikelnummer}
\institute{Hochschule Hamm-Lippstadt -- Software-Ergonomie (AIS-B-2-6.10)}
\date{23. Juni 2026}

\begin{document}

%% ── Titelfolie ──────────────────────────────────────────────────────
\begin{frame}
    \titlepage
\end{frame}

%% ── Agenda ──────────────────────────────────────────────────────────
\begin{frame}{Agenda}
    \tableofcontents
\end{frame}

%% ── 1. Einleitung ───────────────────────────────────────────────────
\section{Einleitung}

\begin{frame}{Problemstellung}
    \begin{itemize}
        \item Website des \textbf{Immersive Reality Lab} (HSHL) mit wachsender Informationsfülle
        \item \textbf{Problem:} Keine interne Suchfunktion, heterogene Seitenstruktur
        \item \textbf{Idee:} Conversational UI -- Fragen in natürlicher Sprache
    \end{itemize}
    \vspace{0.5cm}
    \begin{block}{Ziel}
        RAG-basierter Chatbot \textbf{Ergonomia}, der Fragen zum Lab beantwortet -- \alert{nutzerzentriert gestaltet} nach etablierten Ergonomie-Frameworks.
    \end{block}
\end{frame}

%% ── 2. User Research ──────────────────────────────────────────────
\section{User Research}

\begin{frame}{Personas}
    \begin{tabularx}{\textwidth}{X X X}
        \toprule
        \textbf{Studierende/r} & \textbf{Forschende/r} & \textbf{Industriepartner/in} \\
        \midrule
        Bachelorarbeitsthema & Publikationen suchen & Technologietransfer \\
        Geräteausleihe & Kooperation anbahnen & Ansprechpartner finden \\
        Lehrveranstaltungen & Forschungsgebiete & Gemeinsame Projekte \\
        \bottomrule
    \end{tabularx}
    \vspace{0.5cm}
    \begin{exampleblock}{Top-Findings}
        80~\% würden Chatbot einer Navigation vorziehen. Hauptfrust: unübersichtliche Menüs.
    \end{exampleblock}
\end{frame}

%% ── 3. Theorie ─────────────────────────────────────────────────────
\section{Theoretische Fundierung}

\begin{frame}{Evaluations-Frameworks}
    \begin{itemize}
        \item \textbf{Amershi et al. (2019):} 18 Guidelines for Human-AI Interaction
        \item \textbf{Nielsen (1994):} 10 Usability Heuristiken (adaptiert für Conversational UI)
        \item \textbf{DIN EN ISO 9241-110:} Interaktionsprinzipien
        \item \textbf{SUS:} System Usability Scale (Benchmark: 68)
    \end{itemize}
\end{frame}

%% ── 4. Architektur ─────────────────────────────────────────────────
\section{Architektur \& Design}

\begin{frame}{Systemarchitektur}
    \begin{enumerate}
        \item \textbf{Crawling:} Python + BeautifulSoup scraped Lab-Website
        \item \textbf{Chunking:} 1000 Zeichen, 200 Overlap
        \item \textbf{Embedding:} multilingual-mpnet-base-v2 $\rightarrow$ ChromaDB
        \item \textbf{Retrieval:} Semantische Suche + Keyword-Fallback
        \item \textbf{LLM:} Ollama + Qwen 2.5 (lokal, datenschutzfreundlich)
        \item \textbf{UI:} Next.js 15 + TypeScript + Tailwind
    \end{enumerate}
\end{frame}

\begin{frame}{Ergonomische Designentscheidungen}
    \begin{tabularx}{\textwidth}{lX}
        \toprule
        \textbf{Prinzip} & \textbf{Umsetzung} \\
        \midrule
        Aufgabenangemessenheit & Kurze, präzise Antworten durch RAG \\
        Selbstbeschreibungsfähigkeit & Onboarding-Nachricht + Capability-Response \\
        Fehlertoleranz & ``Weiß ich nicht''-Fallback mit Verweis \\
        Sichtbarkeit des Systemstatus & Ladeindikator + Streaming \\
        Fehlervermeidung & Scope-Begrenzung: Antwort nur aus Kontext \\
        \bottomrule
    \end{tabularx}
\end{frame}

%% ── 5. Demo ────────────────────────────────────────────────────────
\section{Demo}

\begin{frame}{Live-Demo}
    \centering
    \Large
    \alert{Live-Demo des Chatbots}\\
    \vspace{0.5cm}
    \normalsize
    \url{http://localhost:3000}
\end{frame}

%% ── 6. Evaluation ──────────────────────────────────────────────────
\section{Evaluation}

\begin{frame}{SUS-Ergebnisse}
    \centering
    \begin{tabularx}{0.6\textwidth}{lX}
        \toprule
        \textbf{Proband} & \textbf{SUS-Score} \\
        \midrule
        P1 & 72,5 \\
        P2 & 65,0 \\
        P3 & 77,5 \\
        P4 & 70,0 \\
        P5 & 82,5 \\
        \midrule
        \textbf{Mittelwert} & \textbf{73,5} \\
        \bottomrule
    \end{tabularx}
    \vspace{0.5cm}
    \begin{alertblock}{Ergebnis}
        73,5 $\rightarrow$ über Benchmark (68) $\rightarrow$ Note B (``gut'')
    \end{alertblock}
\end{frame}

\begin{frame}{Amershi-Evaluation (Auszug)}
    \begin{tabularx}{\textwidth}{lXc}
        \toprule
        \textbf{Guideline} & \textbf{Bewertung} \\
        \midrule
        G1: Make clear what the system can do & \checkmark \\
        G4: Show contextually relevant information & \checkmark \\
        G7: Support efficient invocation & \checkmark \\
        G10: Scope services when in doubt & \checkmark \\
        G11: Make clear why the system did what it did & \checkmark \\
        G12: Remember recent interactions & \checkmark \\
        G13: Learn from user behavior & -- \\
        G15: Encourage granular feedback & -- \\
        \bottomrule
    \end{tabularx}
    \vspace{0.3cm}
    \small Kern-Guidelines erfüllt; fortgeschrittene Features außerhalb des Projekt-Scope.
\end{frame}

%% ── 7. Fazit ───────────────────────────────────────────────────────
\section{Fazit}

\begin{frame}{Fazit \& Ausblick}
    \begin{itemize}
        \item Ergonomia ist \textbf{kein} Chatbot-Projekt mit UX-Verpackung -- sondern ein \textbf{UX-Projekt}, das einen Chatbot als Medium nutzt
        \item SUS 73,5 bestätigt gute Gebrauchstauglichkeit
    \end{itemize}
    \vspace{0.3cm}
    \begin{block}{Verbesserungspotenzial}
        \begin{itemize}
            \item Vorschlagsfragen im Frontend
            \item Multi-LLM-Strategie (OpenAI/Anthropic als Fallback)
            \item Token-basiertes Chunking
            \item Multi-Session-Management
        \end{itemize}
    \end{block}
\end{frame}

\begin{frame}{Vielen Dank!}
    \centering
    \vspace{1cm}
    \Huge Fragen?
    \vspace{1cm}
    \normalsize
    \url{https://immersive-reality-lab.de}
\end{frame}

\end{document}
